\chapter{Pendahuluan}\label{Pendahuluan}

\pagestyle{plain}

\section{Latar Belakang}

Bab pendahuluan memuat latar belakang atau alasan kuat dilakukannya penelitian, tujuan, dan hipotesis jika ada. Di dalam pendahuluan dijelaskan pula perumusan atau pendekatan penyelesaian masalah dan alasan pemilihan metode yang digunakan. Bergantung pada proses perumusan masalah penelitian, bagian Kerangka Pikir dan Hipotesis dapat ditulis di sini, tidak ditulis dalam bab tersendiri.
Paparan tidak berbelit-belit atau dimulai dengan latar belakang yang terlalu umum. Pernyataan mengenai apa yang diteliti dan apa yang diharapkannya diawali dengan pemikiran logis. \\
 Contoh : \\
	Pariwisata merupakan sektor yang memberikan ciri khas dan daya tarik di suatu wilayah melalui keindahan alam, keberagaman budaya, dan peninggalan sejarah \cite{rahma2020potensi}. 
\section{Rumusan Masalah}
Dari latar belakang dan kerangka pikir, masalah yang diteliti dapat dirumuskan. Masalah yang dirumuskan harus jelas dan fokus pada kata kunci utama yang unik. Dalam merumuskan masalah, deskripsi lokasi studi terutama keunikannya sudah termasuk dalam dalam pertimbangan. Untuk memperjelas perumusan masalah, dapat juga dibuat beberapa pertanyaan yang hendak dijawab dalam penelitian itu. Dalam uraian harus tercakup pendekatan yang digunakan dalam perumusan masalah. \\
Contoh : \\
	Adapun rumusan masalah penelitian ini, sebagai berikut:
\begin{enumerate}
	\item Bagaimana model ARIMA yang dapat dibentuk dari data jumlah wisatawan mancanegara?
	
\end{enumerate}


\section{Tujuan}
Pernyataan tujuan penelitian ialah pernyataan singkat dan jelas tentang tujuan yang akan dicapai sebagai upaya pemecahan masalah maupun memahami gejala (fenomena) yang dijelaskan dalam latar belakang. Gunakan kata kerja yang hasilnya dapat diukur. Bila ada atau memungkinkan, dapat ditulis manfaat atau kegunaan hasil penelitian bagi kepentingan pengembangan ipteks, pertimbangan dalam mengambil kebijakan, kepentingan profesi maupun masyarakat pada umumnya. \\
Contoh : \\
Tujuan dilakukan penelitian ini untuk menganalisis model 